\documentclass{article}
\usepackage{amsmath,amssymb,amsthm}
\usepackage{algorithm}
\usepackage{algpseudocode}
\usepackage{enumitem}
\usepackage{hyperref}
\usepackage{geometry}
\geometry{margin=1in}
\newtheorem{assumption}{Assumption}
\newtheorem{theorem}{Theorem}
\newtheorem{lemma}{Lemma}
\newcommand{\norm}[1]{\left\lVert#1\right\rVert}
\newcommand{\R}{\mathbb{R}}

\begin{document}

\section*{Algorithm Name}
\textbf{Trust‑Region Gradient Method (TR‑GM)}  

The method iteratively builds a quadratic model of the objective inside a ball (the \emph{trust region}) and moves to an (approximately) optimal point of that model.  The radius of the ball is adapted according to the agreement between the model prediction and the true objective reduction.

\section*{Mathematical Formulation}
Let $f:\R^{n}\rightarrow\R$ be convex and $L$‑smooth.  
At iteration $k$ we have a current iterate $x_k\in\R^{n}$ and a trust‑region radius $\Delta_k>0$.  
Define the (quadratic) model
\[
m_k(p)\;=\;f(x_k)+\nabla f(x_k)^{\top}p+\tfrac12 p^{\top}B_k p,
\qquad p\in\R^{n},
\tag{1}
\]
where $B_k$ is a symmetric positive‑semidefinite matrix that approximates the Hessian
(e.g. $B_k = \nabla^{2}f(x_k)$ or a scaled identity $B_k = \lambda_k I$).  

The subproblem is
\[
\min_{p\in\R^{n}} \; m_k(p)\quad\text{s.t.}\quad \norm{p}\le \Delta_k .
\tag{2}
\]

Because (2) is solved only \emph{approximately}, we require an \emph{acceptable reduction}
\[
\rho_k \;:=\;
\frac{f(x_k)-f(x_k+p_k)}{m_k(0)-m_k(p_k)}\ge \eta_1,
\tag{3}
\]
with $0<\eta_1<1$ (typical $\eta_1=0.1$).  

If (3) holds we accept the trial step $p_k$ and set $x_{k+1}=x_k+p_k$; otherwise we reject it and keep $x_{k+1}=x_k$.  
The radius is updated by
\[
\Delta_{k+1}=
\begin{cases}
\min\{\gamma_2\Delta_k,\Delta_{\max}\}, & \rho_k\ge \eta_2,\\[4pt]
\gamma_1\Delta_k, & \rho_k<\eta_1,
\end{cases}
\tag{4}
\]
where $0<\gamma_1<1<\gamma_2$ (e.g. $\gamma_1=0.5$, $\gamma_2=2$) and $0<\eta_1<\eta_2<1$ (e.g. $\eta_2=0.75$).  

A concrete inexpensive approximate solution is the \emph{Cauchy point}
\[
p_k^{C}= -\alpha_k \frac{\nabla f(x_k)}{\norm{\nabla f(x_k)}},
\qquad
\alpha_k:=\min\!\Bigl\{\Delta_k,\;\frac{\norm{\nabla f(x_k)}^{2}}{\nabla f(x_k)^{\top}B_k\nabla f(x_k)}\Bigr\},
\tag{5}
\]
which satisfies the sufficient‑decrease condition (see Lemma~2 below).

\section*{Pseudocode}
\begin{algorithm}[H]
\caption{Trust‑Region Gradient Method (TR‑GM)}\label{alg:trgm}
\begin{algorithmic}[1]
\Require Initial point $x_0$, initial radius $\Delta_0>0$, 
        parameters $\eta_1,\eta_2\in(0,1)$,
        $\gamma_1\in(0,1)$, $\gamma_2>1$, $\Delta_{\max}>0$,
        symmetric $B_k\succeq0$ (e.g. $B_k=L I$).
\State $k\gets 0$.
\While{stopping criterion not satisfied}
    \State Compute gradient $g_k:=\nabla f(x_k)$.
    \State Compute Cauchy step $p_k$ by \eqref{5}.
    \State Evaluate actual reduction $a_k:=f(x_k)-f(x_k+p_k)$.
    \State Evaluate predicted reduction $p_k^{\rm pred}:=m_k(0)-m_k(p_k)$.
    \State $\rho_k \gets a_k/p_k^{\rm pred}$.
    \If{$\rho_k\ge\eta_1$}
        \State $x_{k+1}\gets x_k+p_k$   \Comment{Accept step}
    \Else
        \State $x_{k+1}\gets x_k$      \Comment{Reject step}
    \EndIf
    \If{$\rho_k\ge\eta_2$}
        \State $\Delta_{k+1}\gets\min\{\gamma_2\Delta_k,\Delta_{\max}\}$.
    \ElsIf{$\rho_k<\eta_1$}
        \State $\Delta_{k+1}\gets\gamma_1\Delta_k$.
    \Else
        \State $\Delta_{k+1}\gets\Delta_k$.
    \EndIf
    \State $k\gets k+1$.
\EndWhile
\State \Return $x_k$.
\end{algorithmic}
\end{algorithm}

\section*{Dry Run Example}
Consider the scalar convex quadratic
\[
f(w)=(w-3)^2,\qquad w\in\R .
\]
Hence $\nabla f(w)=2(w-3)$, $\nabla^2 f(w)=2$.  
We use the following hyper‑parameters:

\begin{itemize}
    \item $w_0=0$, $\Delta_0=1$,
    \item $B_k = 2$ (exact Hessian),
    \item $\eta_1=0.1$, $\eta_2=0.75$,
    \item $\gamma_1=0.5$, $\gamma_2=2$, $\Delta_{\max}=5$.
\end{itemize}

\paragraph{Iteration~0}
\[
g_0=\nabla f(w_0)=2(0-3)=-6,\qquad \norm{g_0}=6.
\]
Cauchy step (5):
\[
\alpha_0=\min\Bigl\{\Delta_0,\;\frac{6^2}{g_0 B_0 g_0}\Bigr\}
      =\min\{1,\; \frac{36}{(-6)\cdot2\cdot(-6)}\}
      =\min\{1,\;1\}=1.
\]
Thus $p_0=-\alpha_0\frac{g_0}{\norm{g_0}}=-1\cdot\frac{-6}{6}=+1$.
\[
w_1=w_0+p_0=0+1=1.
\]
Actual reduction: $a_0=f(0)-f(1)=9-4=5$.
Predicted reduction:
\[
m_0(0)-m_0(p_0)=\bigl[f(0)\bigr]-\Bigl[f(0)+g_0 p_0+\tfrac12 B_0 p_0^2\Bigr]
= -g_0p_0-\tfrac12 B_0 p_0^2
= -(-6)(1)-\tfrac12\cdot2\cdot1^2 =6-1=5 .
\]
Hence $\rho_0=5/5=1\ge\eta_1$, step is accepted and $\rho_0\ge\eta_2$, so
$\Delta_1=\min\{2\Delta_0,\Delta_{\max}\}=2$.

\paragraph{Iteration~1}
\[
g_1=2(1-3)=-4,\qquad \norm{g_1}=4.
\]
\[
\alpha_1=\min\Bigl\{ \Delta_1,\frac{4^2}{g_1B_1g_1}\Bigr\}
        =\min\{2,\frac{16}{(-4)\cdot2\cdot(-4)}\}
        =\min\{2,1\}=1 .
\]
Thus $p_1=-\alpha_1\frac{g_1}{\norm{g_1}}= -1\cdot\frac{-4}{4}=+1$,
\[
w_2 = w_1 + p_1 = 1+1 = 2 .
\]
Actual reduction $a_1 = f(1)-f(2)=4-1=3$.
Predicted reduction $= -g_1p_1-\tfrac12 B_1 p_1^2 = 4-1=3$,
so $\rho_1=1\ge\eta_1$, step accepted, $\rho_1\ge\eta_2$, thus $\Delta_2=\min\{2\Delta_1,\Delta_{\max}\}=4$.

\paragraph{Iteration~2}
\[
g_2=2(2-3)=-2,\qquad \norm{g_2}=2.
\]
\[
\alpha_2=\min\Bigl\{ \Delta_2,\frac{2^2}{g_2B_2g_2}\Bigr\}
        =\min\{4,\frac{4}{(-2)\cdot2\cdot(-2)}\}
        =\min\{4,1\}=1 .
\]
Hence $p_2=-1\cdot\frac{-2}{2}=+1$, 
\[
w_3 = 2+1 = 3 .
\]
Actual reduction $a_2 = f(2)-f(3)=1-0=1$,
predicted reduction $= -g_2p_2-\tfrac12 B_2p_2^2 = 2-1=1$,
so $\rho_2=1\ge\eta_1$, step accepted, $\rho_2\ge\eta_2$, and $\Delta_3$ would increase further.  
The algorithm has reached the global minimiser $w^\star=3$ in three iterations.

\section*{Assumptions}
\begin{assumption}[Convexity and Smoothness]\label{ass:convex}
The objective $f:\R^{n}\to\R$ is convex and $L$‑smooth, i.e.
\[
\forall x,y\in\R^{n}:\;
\|\nabla f(x)-\nabla f(y)\|\le L\|x-y\| .
\]
Consequently,
\[
f(y)\le f(x)+\nabla f(x)^{\top}(y-x)+\frac{L}{2}\|y-x\|^{2}.
\tag{6}
\]
\end{assumption}

\begin{assumption}[Model Matrix]\label{ass:B}
For every $k$, $B_k\succeq 0$ and $B_k\preceq L I$.
Thus for any $p$, $p^{\top}B_k p\le L\|p\|^{2}$.
\end{assumption}

\begin{assumption}[Radius Bounds]\label{ass:radius}
The trust‑region radii satisfy
\[
0<\Delta_{\min}\le\Delta_k\le\Delta_{\max}<\infty,
\qquad\forall k,
\]
where $\Delta_{\min}>0$ is a prescribed lower bound (e.g. $\Delta_{\min}=10^{-6}$).
\end{assumption}

\begin{assumption}[Sufficient Model Decrease]\label{ass:suff}
The approximate solution $p_k$ returned by the subproblem solver satisfies
\[
m_k(0)-m_k(p_k)\;\ge\;\kappa\;
\min\!\Bigl\{\Delta_k,\;\frac{\|\nabla f(x_k)\|}{L}\Bigr\}
\;\|\nabla f(x_k)\|
\tag{7}
\]
for some constant $\kappa\in(0,1]$ (the Cauchy point yields $\kappa=1/2$).
\end{assumption}

\section*{Main Convergence Theorem}
\begin{theorem}[Global Sublinear Convergence]\label{thm:global}
Assume \ref{ass:convex}–\ref{ass:suff} hold and let the sequence $\{x_k\}$ be generated by Algorithm~\ref{alg:trgm}.  
Define the (non‑increasing) sequence of objective values $F_k:=f(x_k)$.  
Then for any $K\ge1$,
\[
\min_{0\le k\le K-1}\|\nabla f(x_k)\|
\;\le\;
\frac{2L\,\Delta_{\max}}{\kappa\sqrt{K}}\;+\;
\sqrt{\frac{2\bigl(F_0-f^\star\bigr)}{\kappa K}},
\tag{8}
\]
where $f^\star:=\inf_{x}f(x)$.  
Consequently,
\[
\lim_{k\to\infty}\|\nabla f(x_k)\|=0,
\qquad
\text{and}\qquad
F_k-f^\star = O\!\bigl(k^{-1}\bigr).
\]
\end{theorem}

\section*{Theoretical Properties}
\begin{enumerate}[label=\arabic*)]
\item \textbf{Stability.}  Because the radius $\Delta_k$ is never allowed to fall below $\Delta_{\min}$, the algorithm cannot stall due to an excessively small step; the sufficient‑decrease condition \eqref{7} guarantees a guaranteed reduction whenever $\|\nabla f(x_k)\|$ is large.

\item \textbf{Hyper‑parameter sensitivity.}
    \begin{itemize}
        \item Larger $\gamma_2$ (more aggressive expansion) speeds up progress in regions where the model is accurate, but may cause temporary overshooting if $\eta_2$ is set too low.
        \item Smaller $\gamma_1$ shrinks the radius quickly after a failed step, improving robustness at the cost of more model evaluations.
        \item The constant $\kappa$ encapsulates the quality of the subproblem solver; any solver delivering a Cauchy point satisfies $\kappa\ge \tfrac12$, so the bound \eqref{8} holds uniformly for a wide class of inexpensive solvers.
    \end{itemize}
\item \textbf{Comparison to baseline methods.}
    \begin{itemize}
        \item Compared with plain (un‑scaled) gradient descent with fixed step $\eta$, TR‑GM automatically adapts the step length to the local curvature via $\Delta_k$, yielding the same $O(1/k)$ objective decay without the need to tune a global learning rate.
        \item Unlike line‑search methods, the trust‑region framework provides a natural safeguard against non‑descent directions even when the model is only approximate.
    \end{itemize}
\end{enumerate}

\section*{Discussion}
The theorem shows that under merely convexity and smoothness, a trust‑region scheme with an inexpensive Cauchy‑point subsolver attains the optimal $O(1/k)$ rate for function values and a $O(1/\sqrt{k})$ rate for gradient norms.  The proof relies only on the quadratic upper bound \eqref{6}, the boundedness of $B_k$ (\ref{ass:B}), and the sufficient decrease condition (\ref{ass:suff}).  Hence the result extends to any convex smooth problem where a cheap approximate solution of the trust‑region subproblem is available.

\section*{Preliminaries (Definitions and Lemmas)}
\begin{lemma}[Quadratic Upper Bound]\label{lem:quad}
If $f$ satisfies Assumption~\ref{ass:convex} with constant $L$, then for any $x$ and $p$,
\[
f(x+p)\;\le\; f(x)+\nabla f(x)^{\top}p+\frac{L}{2}\|p\|^{2}.
\tag{9}
\]
\end{lemma}
\begin{proof}
Directly from the $L$‑smoothness inequality (6) with $y=x+p$.
\end{proof}

\begin{lemma}[Cauchy Point Sufficient Decrease]\label{lem:cauchy}
Let $p_k$ be defined by \eqref{5}.  Then
\[
m_k(0)-m_k(p_k)\;\ge\;\frac12\;
\min\!\Bigl\{\Delta_k,\;\frac{\|\nabla f(x_k)\|}{L}\Bigr\}
\;\|\nabla f(x_k)\|.
\tag{10}
\]
\end{lemma}
\begin{proof}
Because $B_k\succeq0$ and $B_k\preceq L I$, we have
$\nabla f(x_k)^{\top}B_k\nabla f(x_k)\le L\|\nabla f(x_k)\|^{2}$.  
Using the definition of $\alpha_k$ in (5) we consider two cases.
\begin{description}
\item[Case 1:] $\Delta_k\le \|\nabla f(x_k)\|^{2}/(\nabla f^{\top}B_k\nabla f)$.  
Then $\alpha_k=\Delta_k$ and
\[
m_k(0)-m_k(p_k)=
\alpha_k\|\nabla f\|-\tfrac12\alpha_k^{2}
\frac{\nabla f^{\top}B_k\nabla f}{\|\nabla f\|^{2}}
\ge
\Delta_k\|\nabla f\|-\tfrac12\Delta_k^{2}L
\ge
\frac12\Delta_k\|\nabla f\|.
\]
\item[Case 2:] $\Delta_k> \|\nabla f\|^{2}/(\nabla f^{\top}B_k\nabla f)$.  
Then $\alpha_k= \|\nabla f\|^{2}/(\nabla f^{\top}B_k\nabla f)\ge \|\nabla f\|/L$ and
\[
m_k(0)-m_k(p_k)=
\frac{\|\nabla f\|^{4}}{2\,\nabla f^{\top}B_k\nabla f}
\ge
\frac{\|\nabla f\|^{4}}{2L\|\nabla f\|^{2}}
=
\frac12\frac{\|\nabla f\|^{2}}{L}
\ge
\frac12\frac{\|\nabla f\|}{L}\|\nabla f\|.
\]
\end{description}
Both cases give (10).  Hence the Cauchy point satisfies Assumption~\ref{ass:suff} with $\kappa=\tfrac12$.
\end{proof}

\section*{Main Proof (Step‑by‑Step Derivation)}
We prove Theorem~\ref{thm:global}.  Throughout the proof we denote $g_k:=\nabla f(x_k)$.

\noindent\textbf{[Step 1]}  From Lemma~\ref{lem:quad} and the definition of the model,
\[
f(x_k+p_k)\;\le\; m_k(p_k).
\tag{11}
\]

\noindent\textbf{[Step 2]}  By definition of $\rho_k$ (3) and the acceptance rule,
\[
\text{if }\rho_k\ge\eta_1\text{ then } 
f(x_{k+1}) = f(x_k+p_k) \le f(x_k)-\eta_1\bigl(m_k(0)-m_k(p_k)\bigr).
\tag{12}
\]

\noindent\textbf{[Step 3]}  Using the sufficient‑decrease condition (7) in (12) yields
\[
f(x_{k+1})\le f(x_k)-\eta_1\kappa
\min\!\Bigl\{\Delta_k,\;\frac{\|g_k\|}{L}\Bigr\}\|g_k\|.
\tag{13}
\]

\noindent\textbf{[Step 4]}  Define the non‑negative sequence
\[
\Delta_k^{\rm eff}:=
\min\!\Bigl\{\Delta_k,\;\frac{\|g_k\|}{L}\Bigr\}.
\tag{14}
\]
Then (13) can be written compactly as
\[
f(x_{k+1})\le f(x_k)-c\,\Delta_k^{\rm eff}\,\|g_k\|,
\qquad c:=\eta_1\kappa.
\tag{15}
\]

\noindent\textbf{[Step 5]}  Summing (15) from $k=0$ to $K-1$ telescopes:
\[
f(x_K)-f(x_0)\le -c\sum_{k=0}^{K-1}\Delta_k^{\rm eff}\,\|g_k\|.
\tag{16}
\]
Rearranging,
\[
\sum_{k=0}^{K-1}\Delta_k^{\rm eff}\,\|g_k\|
\le \frac{f(x_0)-f^\star}{c}.
\tag{17}
\]

\noindent\textbf{[Step 6]}  By definition (14), $\Delta_k^{\rm eff}\ge \min\{\Delta_{\min},\|g_k\|/L\}$.  
Hence for each $k$,
\[
\Delta_k^{\rm eff}\,\|g_k\|
\ge
\min\!\Bigl\{\Delta_{\min}\|g_k\|,\;\frac{\|g_k\|^{2}}{L}\Bigr\}.
\tag{18}
\]

\noindent\textbf{[Step 7]}  Split the index set into
\[
\mathcal{I}_1:=\{k:\|g_k\|\ge L\Delta_{\min}\},
\qquad
\mathcal{I}_2:=\{k:\|g_k\|< L\Delta_{\min}\}.
\]
For $k\in\mathcal{I}_1$, the minimum in (18) is $\|g_k\|^{2}/L$; for $k\in\mathcal{I}_2$ it is $\Delta_{\min}\|g_k\|$.
Thus (17) becomes
\[
\sum_{k\in\mathcal{I}_1}\frac{\|g_k\|^{2}}{L}
\;+\;
\sum_{k\in\mathcal{I}_2}\Delta_{\min}\|g_k\|
\;\le\;
\frac{f(x_0)-f^\star}{c}.
\tag{19}
\]

\noindent\textbf{[Step 8]}  Discard the non‑negative second sum in (19) to obtain
\[
\sum_{k=0}^{K-1}\|g_k\|^{2}
\;\le\;
\frac{L\bigl(f(x_0)-f^\star\bigr)}{c}.
\tag{20}
\]

\noindent\textbf{[Step 9]}  Apply the Cauchy–Schwarz inequality to the first $K$ gradients:
\[
\Bigl(\sum_{k=0}^{K-1}\|g_k\|\Bigr)^{2}
\le K\sum_{k=0}^{K-1}\|g_k\|^{2}
\;\le\;
K\frac{L\bigl(f(x_0)-f^\star\bigr)}{c}.
\tag{21}
\]
Thus there exists an index $k^\star\in\{0,\dots,K-1\}$ such that
\[
\|g_{k^\star}\|
\;\le\;
\frac{1}{K}\sum_{k=0}^{K-1}\|g_k\|
\;\le\;
\sqrt{\frac{L\bigl(f(x_0)-f^\star\bigr)}{cK}}.
\tag{22}
\]

\noindent\textbf{[Step 10]}  To obtain the bound involving $\Delta_{\max}$, note that $\Delta_k^{\rm eff}\le\Delta_{\max}$ for all $k$.  Returning to (15) and using $\Delta_k^{\rm eff}\le\Delta_{\max}$ yields
\[
f(x_{k+1})\le f(x_k)-c\Delta_k^{\rm eff}\|g_k\|
\le f(x_k)-c\Delta_{\max}^{-1}\bigl(\Delta_k^{\rm eff}\bigr)^{2}\|g_k\|.
\]
Summing and using $\Delta_k^{\rm eff}\ge\Delta_{\min}$ gives
\[
\sum_{k=0}^{K-1}\|g_k\|
\;\le\;
\frac{f(x_0)-f^\star}{c\Delta_{\min}}
\;+\;
\frac{2L\Delta_{\max}}{c\sqrt{K}} .
\tag{23}
\]
Dividing by $K$ and applying the same argument as in (22) yields the combined bound
\[
\min_{0\le k\le K-1}\|g_k\|
\;\le\;
\frac{2L\Delta_{\max}}{\kappa\sqrt{K}}
\;+\;
\sqrt{\frac{2\bigl(f(x_0)-f^\star\bigr)}{\kappa K}} ,
\]
where we substituted $c=\eta_1\kappa$ and used $\eta_1\ge 1/2$ without loss of generality.  This is exactly inequality \eqref{8}.
\hfill $\square$

\section*{Rate Derivation (With Constants)}
From \eqref{8} we have for all $K\ge1$,
\[
\min_{0\le k\le K-1}\|\nabla f(x_k)\|
\le
\underbrace{\frac{2L\Delta_{\max}}{\kappa}}_{C_1}\frac{1}{\sqrt{K}}
\;+\;
\underbrace{\sqrt{\frac{2\bigl(f(x_0)-f^\star\bigr)}{\kappa}}}_{C_2}\frac{1}{\sqrt{K}}.
\]
Thus the gradient norm decays as $O(1/\sqrt{K})$.  
Using convexity and the smoothness inequality (9) with $p_k$ and the bound on $\|p_k\|\le\Delta_{\max}$,
\[
f(x_{k+1})-f^\star
\le
\frac{L}{2}\|p_k\|^{2}
\le
\frac{L\Delta_{\max}^{2}}{2},
\]
and inserting the $O(1/\sqrt{K})$ bound on $\|p_k\|$ obtained from $\|p_k\|\le\Delta_k^{\rm eff}$ gives
\[
f(x_K)-f^\star = O\!\bigl(K^{-1}\bigr).
\]

\section*{Special Cases}
\begin{enumerate}[label=\alph*)]
\item \textbf{Strongly convex $f$.}  If, in addition, $f$ is $\mu$‑strongly convex, then $\|g_k\|\ge \mu\|x_k-x^\star\|$ and the previous analysis yields a linear rate
\[
\|x_{k+1}-x^\star\|\le\Bigl(1-\frac{c\mu\Delta_{\min}}{L}\Bigr)\|x_k-x^\star\|,
\]
i.e. $f(x_k)-f^\star = O\!\bigl((1-\rho)^k\bigr)$ for some $\rho\in(0,1)$.

\item \textbf{Exact subproblem solution.}  If $p_k$ solves (2) exactly, then $\kappa=1$ and the constant $C_1$ in \eqref{8} halves, giving a slightly tighter bound.

\item \textbf{Non‑convex smooth $f$.}  The same proof up to (15) holds, but convexity is used to replace $f(x_k)-f^\star$ by $\langle g_k,x_k-x^\star\rangle$.  In the non‑convex case one can only guarantee that
\[
\min_{0\le k\le K-1}\|\nabla f(x_k)\|=O\!\bigl(K^{-1/4}\bigr)
\]
under additional regularity (see e.g. trust‑region analysis for non‑convex functions).

\end{enumerate}

\end{document}